\chapter{Flexbox in integers}
\label{ch:layout}

Given an element tree, \maya{} must decide where every text run, every
box, every border goes --- in cells, not pixels. The algorithm it uses
is a tailored re-implementation of CSS Flexbox. This chapter explains
flexbox from scratch (no web background assumed) and then walks
\maya{}'s implementation.

\begin{aside}
Layout is the unglamorous middle of every UI library. Nobody talks
about it; the algorithm runs hundreds of times a second, makes
thousands of small decisions, and disappears the moment it gets the
rectangle right. Bad layout is instantly noticeable; good layout is
invisible. This chapter is a guided tour of \maya{}'s ten kilobytes
of layout code, because once you understand it you can stop guessing
at why a panel is one cell off and start reading the cause out of
the configuration.
\end{aside}

\section{The problem layout solves}

You have a rectangular area --- the terminal --- and a tree of
elements each with constraints (\emph{this one wants to be 20 cells
wide; this row should pack its children with a 1-cell gap; this panel
must grow to fill leftover space}). Layout is the function that
assigns concrete \code{(x, y, w, h)} rectangles to each node such that
all the constraints are satisfied.

There are many possible algorithms (table layout, constraint solving,
grid layout). Flexbox is the most appropriate for terminal UIs:
predictable, single-pass, no surprises, no solver.

\subsection{Why not the obvious algorithms?}

A brief tour of the alternatives \maya{} did not choose:

\paragraph{Table layout (HTML \code{<table>}).} A rectangular grid
of rows and columns. Each cell's size is determined by the largest
content in its row and column. Beautiful for tabular data; awkward
for anything else --- two unrelated cells affecting each other's
size is a constant surprise. Used in early web design; abandoned
as a general-purpose layout in the mid-2000s.

\paragraph{Constraint solving (Cassowary, Auto Layout).} Each
position is a variable; each layout rule is a linear constraint.
A simplex-like solver finds an assignment. Powerful (it can express
``these two are always equal'', ``this is between 10 and 30
cells''), but expensive and prone to producing surprising answers
when constraints conflict. Apple's Auto Layout works this way.

\paragraph{Grid (CSS Grid).} A two-dimensional placement model:
you declare a grid of named rows and columns, then place children
in specific cells. Excellent for page-level layout; overkill for
the nested-row-and-column shape most TUIs take.

\paragraph{Flow / line-box (HTML inline text).} Children are
packed left-to-right, top-to-bottom, like words in a paragraph.
Good for prose; bad for everything else.

Flexbox is the simplest model that handles the common cases ---
rows of widgets, columns of panels, a sidebar plus content, a
status bar pinned to the bottom --- without ceremony. It is also
linear-time and predictable, both crucial for a sixty-frames-per-second
renderer.

\begin{funfact}
Flexbox went through several names and standards before it stuck.
The original CSS3 ``flexible box'' module was proposed in 2009;
the syntax changed twice (the \code{display: box} from 2009, then
\code{display: flexbox} in 2011, then \code{display: flex} in 2012)
before browsers settled. Mosaic's table-based layouts and
Netscape's \code{<center>} tag ran the web for a decade because
nothing nicer was standardised. The model \maya{} uses --- main
axis, grow, shrink, basis, gap --- is the 2012 design, late enough
that the rough edges had been filed off, simple enough to fit in
800 lines of \cpp{}.
\end{funfact}

\section{Flexbox in one paragraph}

Flexbox arranges children along a \textbf{main axis} (row or column).
Children may \emph{grow} to fill leftover space and \emph{shrink} when
oversized; they may \emph{wrap} onto new lines when they don't fit;
they may be \emph{aligned} along the cross axis and \emph{justified}
along the main axis. Each container picks an axis (row or column);
each child carries a few numbers (basis, grow, shrink, gap, alignment
override).

\section{Vocabulary}

A worked example will be clearest. Take a horizontal row of three
children inside a 30-cell-wide box:

\begin{lstlisting}
h(
    panel_a,   // wants 10 cells
    panel_b,   // grow=1
    panel_c    // wants 8 cells
) | gap_<1>
\end{lstlisting}

\maya{} reasons about it like this:

\begin{enumerate}[leftmargin=*]
  \item \textbf{Sum the basis sizes.} A=10, B=0 (no fixed size,
    grows), C=8. Plus 2 gap cells. Total committed: 20.
  \item \textbf{Leftover space.} 30 - 20 = 10 cells.
  \item \textbf{Distribute by grow.} Only B has \code{grow > 0}, so
    B receives all 10 cells. New size: B = 10.
  \item \textbf{Place.} A at column 0, gap, B at column 11, gap, C
    at column 22.
\end{enumerate}

The cross axis (vertical, in this example) is handled by
\code{Align}: children stretch to the parent height by default, but
they can be \code{Start}, \code{Center}, \code{End}, or
\code{Stretch}.

\section{The flex properties}

A node in \maya{}'s layout engine carries the following style:

\begin{lstlisting}
struct FlexStyle {
    FlexDirection direction = FlexDirection::Column;
    FlexWrap      wrap      = FlexWrap::NoWrap;

    Align   align_items   = Align::Auto;
    Align   align_self    = Align::Auto;
    Justify justify       = Justify::Auto;

    int     gap           = 0;

    int     pad_top    = 0, pad_right = 0,
            pad_bottom = 0, pad_left  = 0;
    int     mar_top    = 0, mar_right = 0,
            mar_bottom = 0, mar_left  = 0;

    int     width       = -1;   // -1 = auto
    int     height      = -1;
    int     min_width   = 0, max_width  = INT_MAX;
    int     min_height  = 0, max_height = INT_MAX;

    int     basis       = -1;   // preferred main-axis size, -1 = auto
    float   grow        = 0.0f;
    float   shrink      = 1.0f;

    Display display     = Display::Flex;
    Overflow overflow   = Overflow::Visible;
    bool    has_border  = false;
};
\end{lstlisting}

These mirror CSS flexbox almost exactly. The differences:

\begin{itemize}[leftmargin=*]
  \item Every quantity is an \emph{integer} number of cells. No
    \code{em}, no \code{px}, no fractional pixels. Sub-cell positions
    don't exist in a terminal.
  \item \code{display: Flex} and \code{display: None} are the only
    options. There is no block, no inline.
  \item Percentages are limited; most sizes are explicit integers.
\end{itemize}

\section{The data structure}

A layout tree is a flat \code{std::vector<LayoutNode>} indexed by
\code{int}. Each node has a list of child indices.

\begin{lstlisting}
struct LayoutNode {
    FlexStyle           style;
    std::vector<int>    children;     // indices into the same vector
    layout::Rect        computed;     // output: x, y, w, h
};
\end{lstlisting}

Why a flat vector and not a pointer tree? Three reasons:

\begin{itemize}[leftmargin=*]
  \item \textbf{Cache locality.} Walking siblings is a stride of
    \code{sizeof(LayoutNode)}; the prefetcher loves it.
  \item \textbf{No allocation in the hot path.} A
    \code{vector<int>} for children is cheaper than a \code{vector
    <unique\_ptr<LayoutNode>>}.
  \item \textbf{Easier reasoning.} Indices cannot dangle; an index
    is either in range or out of range. Pointer trees create
    aliasing puzzles.
\end{itemize}

\begin{idea}
``Make it an index, not a pointer'' is a recurring optimisation in
\maya{}. The same idea drives the style pool (a 16-bit ID into a
table) and the hyperlink interning (a 32-bit ID).
\end{idea}

\section{The algorithm, step by step}

\code{layout::compute(nodes, root, available\_width)} runs in three
phases per box:

\begin{enumerate}[leftmargin=*]
  \item \textbf{Measure.} Walk children recursively and compute their
    intrinsic main-axis size (the basis if specified, else the
    content's measured size).
  \item \textbf{Distribute.} Add up basis sizes plus gaps. Compare
    to the parent's main-axis size. If there is leftover space and
    any child has \code{grow > 0}, distribute the leftover
    proportionally; if there is overflow and any child has
    \code{shrink > 0}, shrink proportionally.
  \item \textbf{Place.} Compute final \code{(x, y, w, h)} for each
    child given the resolved main-axis sizes and the cross-axis
    alignment.
\end{enumerate}

Wrapping (when \code{wrap = Wrap}) re-enters the algorithm per line.
Justify and align use simple integer division: \code{SpaceBetween}
puts the leftover between children, \code{SpaceAround} puts it on
both sides, \code{Center} halves it.

\subsection{Padding, border, margin}

Three rectangles concentric:

\begin{center}
\small
\begin{tabular}{ll}
\textbf{margin} & cells \emph{outside} the box (sibling spacing) \\
\textbf{border} & one cell thick if present (drawn glyphs) \\
\textbf{padding} & cells \emph{inside} the border (content inset) \\
\textbf{content} & where children are laid out \\
\end{tabular}
\end{center}

The available width for children is \code{w - left\_pad - right\_pad -
border\_l - border\_r}. The available height is the same on the
vertical. \maya{}'s renderer paints the border in the one-cell ring
around the padded content.

\section{Worked example: a status panel}

\begin{lstlisting}
auto panel = v(
    t<"Status"> | Bold,
    sep,
    h(t<"CPU:"> | Dim, space, t<"42%"> | Bold) | gap_<1>,
    h(t<"Mem:"> | Dim, space, t<"8.2G"> | Bold) | gap_<1>
) | border_<Round> | pad<1, 2> | gap_<1>;
\end{lstlisting}

Suppose the parent box gives us 24 cells wide. The layout proceeds:

\begin{enumerate}[leftmargin=*]
  \item Outer box: border (1 cell each side) + padding (2 on each
    side) leaves 24 - 2 - 4 = \textbf{18 cells} of inner width.
  \item Inner column has four children: title (1 row), separator (1
    row), CPU line (1 row), Mem line (1 row), with gap 1 between
    them. Total height 4 + 3 = 7 rows of content. Plus the border
    pads (1 + 1) and padding (1 + 1) makes the outer box 11 rows.
  \item Inside each \code{h(...)} line, the children are: a static
    label, a \code{space} (grow=1), and a value text. Total fixed
    width: maybe 4 + 3 + gap 1 = 8 cells. Leftover 10 goes entirely
    to the spacer. The value text is pushed to column 8 + 10 = 18,
    i.e. the right edge.
\end{enumerate}

The renderer then walks the resulting rectangles and writes glyphs.

\section{When layout disagrees with you}

Two situations confuse newcomers:

\begin{itemize}[leftmargin=*]
  \item \textbf{Nothing happens when I set \code{grow}.} The parent
    must have a fixed (or grow-controlled) size for the leftover
    space to be meaningful. A \code{grow=1} child of a parent that
    sizes itself to its content has nothing to grow into.
  \item \textbf{Text wraps inside a column unexpectedly.} A
    \code{TextElement}'s width is its content. If you put a long
    string inside a column, it tries to take its natural width; the
    parent box constrains it to its own width, which causes wrap.
    Set \code{TextWrap::Truncate} or \code{Ellipsize} if that is
    wrong.
\end{itemize}

\section{Where the code lives}

\file{maya/include/maya/layout/yoga.hpp} and
\file{maya/src/layout/yoga.cpp}. Despite the name, this is not a
binding to Facebook's Yoga library --- it is an independent
implementation that follows the same flexbox spec. ``Yoga'' here is a
nickname for the algorithm. The implementation is around 800 lines
and has no external dependencies.

\section{Justify and Align, every value}
\label{sec:layout:justify-align}

The two enums \code{Justify} and \code{Align} control how leftover
space is distributed and how children line up across the cross axis.
Both share a value list: \code{Start}, \code{End}, \code{Center},
\code{Stretch}, \code{SpaceBetween}, \code{SpaceAround},
\code{SpaceEvenly}, \code{Auto}. Worth walking each.

Imagine a row 30 cells wide, holding three children each 4 cells
wide. Leftover space: $30 - 12 = 18$ cells.

\paragraph{\code{Start} (default).} Pack children to the start.
No space inserted between them; all 18 leftover cells go to the
end.

\begin{terminalbox}
AAAA BBBB CCCC..................
\end{terminalbox}

\paragraph{\code{End}.} Pack children to the end. All leftover
space at the start.

\begin{terminalbox}
..................AAAA BBBB CCCC
\end{terminalbox}

\paragraph{\code{Center}.} Half the leftover before, half after.
With odd leftover (rare in cells but possible), the extra cell
goes to the start.

\begin{terminalbox}
.........AAAA BBBB CCCC.........
\end{terminalbox}

\paragraph{\code{SpaceBetween}.} Leftover divided into
$n - 1$ equal gaps between children. First child at start, last at
end. With 18 cells divided into 2 gaps: 9 each.

\begin{terminalbox}
AAAA.........BBBB.........CCCC
\end{terminalbox}

\paragraph{\code{SpaceAround}.} Leftover divided into $n$ equal
segments, half a segment on each end, full segments between.
With 18 cells over 3 children: 6 per segment, 3 on each end.

\begin{terminalbox}
...AAAA......BBBB......CCCC...
\end{terminalbox}

\paragraph{\code{SpaceEvenly}.} Leftover divided into $n + 1$
equal gaps. Same gap on each end as between children. With 18
cells over 4 gaps: 4 or 5 each (integer division leaves a remainder).

\begin{terminalbox}
....AAAA....BBBB....CCCC...
\end{terminalbox}

\paragraph{\code{Stretch}.} Cross-axis only. Children are sized to
the parent's cross-axis size. The default for cross-axis alignment
on boxes (so children fill the height of a row by default).

\paragraph{\code{Auto}.} The default. For \code{align\_items}
resolves to \code{Stretch}; for \code{align\_self} means ``inherit
from parent's \code{align\_items}''.

\subsection{Integer division and the remainder}

A wrinkle of cell-based layout: leftover space rarely divides
evenly. With \code{SpaceBetween} of three children in 30 cells
(leftover 18, two gaps), each gap is 9. With four children in 30
cells (leftover 14, three gaps), each gap is 4 and the integer
division leaves 2 cells unaccounted for.

\maya{}'s strategy: the first $k$ gaps get one extra cell, where
$k$ is the remainder. Children later in the list see a slightly
tighter gap. The alternative would be to round each gap and
accept that the total drifts by a cell or two, but accumulating
rounding errors quickly produces visible misalignment.

\begin{idea}
In integer cell-based layout, leftover space rarely divides
evenly. The convention is to give the remainder to the early gaps,
so the off-by-one shows up at one edge rather than accumulating.
This is a recurring theme in pixel-free graphics; the same
problem appears in audio sample rate conversion and frame-rate
remapping.
\end{idea}

\section{Grow and shrink, mechanically}
\label{sec:layout:grow-shrink}

The \code{grow} and \code{shrink} floats are the most subtle part
of flexbox. Worth walking through their full mechanics, with
worked examples.

\subsection{The setup}

Given children with basis sizes $b_1, b_2, \dots, b_n$ and an
available main-axis size $W$:

\begin{enumerate}[leftmargin=*]
  \item Sum the basis sizes and gaps: $B = \sum b_i + \text{gaps}$.
  \item If $B = W$, you are done. Place children at their basis
    sizes.
  \item If $B < W$: there is leftover space, and \textbf{grow}
    distributes it. If no child has \code{grow > 0}, leftover stays
    leftover (handled by \code{justify}).
  \item If $B > W$: children overflow, and \textbf{shrink}
    distributes the negative. If no child has \code{shrink > 0},
    the overflow stays (visible as content escaping the box, or
    clipped depending on \code{overflow}).
\end{enumerate}

\subsection{Distributing leftover by \code{grow}}

Let $G = \sum g_i$ be the total grow factor of all children. If
$G > 0$ and $L = W - B > 0$, each child $i$ with grow $g_i$ receives
an additional $\lfloor L \cdot g_i / G \rfloor$ cells.

A worked example. Three children, basis sizes 10, 5, 5 (total 20),
grow factors 1, 0, 2 (total 3), available 32 cells. Leftover
$L = 12$. Distribute:

\begin{itemize}[leftmargin=*]
  \item Child 1: $\lfloor 12 \cdot 1 / 3 \rfloor = 4$ extra cells.
    Final size 14.
  \item Child 2: grow 0, no extra. Final size 5.
  \item Child 3: $\lfloor 12 \cdot 2 / 3 \rfloor = 8$ extra cells.
    Final size 13.
\end{itemize}

Total: $14 + 5 + 13 = 32$, exactly $W$. The remainders from
division sum to zero in this case; when they don't, the first child
in the list absorbs the difference (so the total comes out exact).

\subsection{Distributing overflow by \code{shrink}}

The symmetric case. If $B > W$ and total shrink $S > 0$, each child
loses $\lfloor (B - W) \cdot s_i / S \rfloor$ cells. The defaults
are important: \code{shrink = 1.0f} means every child shrinks
equally; \code{shrink = 0.0f} means the child refuses to shrink and
overflow goes elsewhere.

A child whose basis is already at its \code{min\_width} cannot
shrink further; the algorithm marks it ineligible and redistributes
over the remaining children. This is one of two places \maya{}'s
flexbox implementation does a second pass.

\subsection{Common patterns}

The spinner of common grow/shrink configurations:

\begin{center}
\small
\begin{tabular}{ll}
\textbf{Effect} & \textbf{Setting} \\
\hline
Fixed size, never resizes        & \code{grow=0, shrink=0, basis=N} \\
Fills leftover space              & \code{grow=1, shrink=1} (defaults) \\
Fills, never overflows            & \code{grow=1, shrink=1, min\_width=0} \\
Fills proportionally (2:1 split)  & two siblings with \code{grow=2} and \code{grow=1} \\
Never shrinks below content       & \code{shrink=0} \\
Doesn't grow, will shrink         & \code{grow=0, shrink=1} \\
\end{tabular}
\end{center}

\code{space} in the DSL is a node with \code{grow=1, basis=0,
shrink=0}: it eagerly absorbs every leftover cell and refuses to be
shrunk if there isn't enough space.

\section{Why integer cells, not floats}
\label{sec:layout:integer-math}

Web flexbox uses floating-point pixels because the visual unit ---
a pixel --- can effectively be subdivided in CSS (sub-pixel
antialiasing, fractional positions, transforms). The terminal has
no such unit: a cell is a cell, you cannot draw half a character.
Layout has to terminate on whole cells.

Using \code{int} for every quantity has consequences worth
being explicit about:

\paragraph{No rounding surprises during composition.} Two layouts
computed independently and then combined have exactly the same
geometry as if computed together. With floats, rounding error
compounds: a panel that is 12.499... cells wide at one level
becomes 12 at the next, then 11.998... in the parent, and so on.
Integer math avoids the question.

\paragraph{Diff is bitwise comparable.} Frame-to-frame layout
diffs (Chapter~\ref{ch:simd}) can compare integer rectangles
with a single \code{==}; with floats the question of ``did this
position change'' becomes ``did this position change by more than
epsilon'', and the epsilon has to be chosen.

\paragraph{Constraints are exact.} ``This box is at least 10 cells
wide and at most 30'' is bit-exact; the solver never needs to
worry about a width that is mathematically 10.0001 cells.

\paragraph{The cost: integer division does not preserve totals.}
The technique from the previous section --- give remainders to the
earliest gaps --- is what we pay for using integers. It is much
less work than the floating-point sorrows we avoid.

\section{The three phases, expanded}
\label{sec:layout:phases}

The summary said ``measure, distribute, place''. The full version
for a single box:

\paragraph{Phase 1: measure.}
For each child, recursively compute its \emph{intrinsic} main-axis
size --- the size it wants to be if given no extra space.
\code{TextElement}'s intrinsic main-axis size is the width of its
content (for a row container) or the height after wrapping (for a
column). \code{BoxElement}'s intrinsic is its own measured size,
computed by recursing into its children.

The result is each child's \textbf{basis} --- either the explicit
\code{basis} field if set, or the measured intrinsic.

\paragraph{Phase 2: distribute.}
Compare $B = \sum b_i + \text{gaps}$ to $W$, the available
main-axis size. Distribute leftover or overflow per
Section~\ref{sec:layout:grow-shrink}. After this phase every
child has a final main-axis size.

\paragraph{Phase 3: place.}
Walk children, assigning each one an $(x, y)$. The cursor advances
by the child's final main-axis size plus the gap; the cross-axis
position is determined by \code{align}. If \code{justify} is
non-trivial, the gap is enlarged or padding inserted at the ends.

Each child's final \code{computed} rectangle is written into the
\code{LayoutNode} struct, and the algorithm recurses into the
child to lay out \emph{its} children, with the child's resolved
size as the new available area.

\subsection{Wrapping}

When \code{wrap = Wrap} and the children don't fit on one line,
phase 2 partitions them into lines greedily: pack until the next
child would overflow, start a new line. Each line is then
independently laid out and placed at increasing cross-axis offsets.

For a wrapping column (rare), the lines are columns and the cross
axis is horizontal. Symmetric.

\section{Debugging layout: the four-step ritual}
\label{sec:layout:debugging}

When something is one cell off, two cells over, or in the wrong
place entirely, the cause is almost always one of four things.
Worth checking each in order before opening a debugger.

\paragraph{1. Check the available size.} The terminal's columns
and rows are passed in at the top of the layout. Print them. A
resize between frames will sometimes lag a frame, and you might be
laying out into the wrong rectangle.

\paragraph{2. Check padding and border on every ancestor.} A box
with padding 1 and a border eats 4 cells of width before its
children see anything. Inside a $24$-column terminal, after a
bordered-padded panel and another bordered-padded inner box,
children get 24 - 4 - 4 = 16 cells. The defect is rarely the layout;
it is usually an extra padding nobody remembered.

\paragraph{3. Print the layout tree.} A small helper:

\begin{lstlisting}
void dump(const LayoutNode& n, int depth = 0) {
    std::string indent(depth * 2, ' ');
    std::println("{}{}: ({},{}) {}x{} basis={} grow={}",
        indent, kind_name(n),
        n.computed.x, n.computed.y,
        n.computed.w, n.computed.h,
        n.style.basis, n.style.grow);
    for (int c : n.children) dump(nodes[c], depth + 1);
}
\end{lstlisting}

The answer is almost always visible in the first three lines of
output.

\paragraph{4. Use the DSL's debug border.} Wrap the suspicious
node in a border (\code{| border\_<Single>}); the corner glyphs
show you where its rectangle actually is. If the border is one
cell off, the layout is one cell off.

\begin{tryit}
Run any \maya{} sample and resize your terminal slowly. The
layout reflows on every \code{SIGWINCH}; you should see boxes
rearrange in cell steps, never half a cell. Try narrowing until
your status bar wraps; the wrap point is where total basis +
gaps exceeds the new available width.
\end{tryit}

\section{Layout recipes}
\label{sec:layout:recipes}

A handful of common shapes and how to spell them in the DSL.

\paragraph{Sidebar + content.}

\begin{lstlisting}
h(
    sidebar  | w_<24>,
    content  | grow_<1>
)
\end{lstlisting}

Sidebar fixed at 24; content takes everything else.

\paragraph{Header + body + footer.}

\begin{lstlisting}
v(
    header,
    body | grow_<1>,
    footer
)
\end{lstlisting}

Header and footer take their natural height; body fills the
leftover vertical space.

\paragraph{Three-pane equal split.}

\begin{lstlisting}
h(
    pane_a | grow_<1>,
    pane_b | grow_<1>,
    pane_c | grow_<1>
) | gap_<1>
\end{lstlisting}

Three equal-width panes with one-cell gaps. Each gets
$(W - 2) / 3$ cells.

\paragraph{Status bar on the right.}

\begin{lstlisting}
h(
    text(left),
    space,                  // spacer (grow=1, basis=0)
    text(right) | Dim
)
\end{lstlisting}

The spacer absorbs all leftover space, pushing the right text to
the edge.

\paragraph{Centred dialog.}

\begin{lstlisting}
h(
    space,
    v(
        space,
        dialog | w_<60>,
        space
    ),
    space
)
\end{lstlisting}

Four spacers form a cross-shaped frame around a 60-column dialog,
centring it horizontally and vertically.

\paragraph{A list with a fixed header.}

\begin{lstlisting}
v(
    header_row | Bold,
    sep,
    map(items, render_row) | grow_<1>
)
\end{lstlisting}

Header takes its row, separator one row, the body takes the rest.
When the list overflows, the \code{| grow\_<1>} on the mapped
content makes its container scrollable (Chapter~\ref{ch:widgets}
covers scroll surfaces).

\section{Performance notes}
\label{sec:layout:perf}

Layout is $O(n)$ in the number of nodes for an unwrapped tree.
With wrap, the worst case is $O(n^2)$ (every child triggers a
re-evaluation of remaining line capacity), but typical trees do
not wrap and the constant factor is small.

Numbers from the maya benchmark suite, on a typical desktop CPU:

\begin{itemize}[leftmargin=*]
  \item A 50-node tree (typical small app): under 5\,\textmu{}s.
  \item A 500-node tree (a long list of items): around
    30\,\textmu{}s.
  \item A 5000-node tree (pathological --- a maximalist
    dashboard): around 300\,\textmu{}s.
\end{itemize}

Layout is rarely the bottleneck. If it shows up in your profiler,
the usual cause is rebuilding too many DSL nodes per frame, not
the solver itself. Flatten where you can; cache \code{constexpr}
subtrees that do not change.

\section{Pitfalls and anti-patterns}
\label{sec:layout:pitfalls}

\begin{warning}
\textbf{Negative space from over-applied padding.} If parent and
child both apply padding, the visible insets compound. Decide which
layer owns the padding and keep it there.
\end{warning}

\begin{warning}
\textbf{A grow child inside a non-growing parent.} The parent sizes
itself to content, so there is no leftover space to distribute, and
the child does not grow. To get a growing layout, every ancestor up
to the root must be growing or explicitly sized.
\end{warning}

\begin{warning}
\textbf{Mixing \code{shrink=0} children with insufficient width.}
If every child refuses to shrink and they collectively exceed the
parent's width, content overflows. The renderer clips; the user
sees text cut off mid-glyph. Allow at least the variable-content
child to shrink.
\end{warning}

\begin{warning}
\textbf{Forgetting that text has intrinsic width.} A
\code{TextElement} in a row sizes to its content width by default.
Long strings push siblings off the edge. Wrap them in a fixed-width
box (\code{| w\_<30>}) or set their wrap mode to
\code{TextWrap::Truncate}.
\end{warning}

\begin{recap}
Layout decides where things go. \maya{} uses an integer-cell flexbox
solver over a flat \code{vector<LayoutNode>} addressed by index.
Three phases: measure, distribute, place. The same model handles
rows, columns, wrapping, alignment, justification, and nested boxes.
\end{recap}

\section*{Workshop}

\begin{enumerate}[leftmargin=*]
  \item Write a row that displays your name on the left and the date
    on the right, with the space between them growing to fill the
    terminal width.
  \item Build a 3-column dashboard: three boxes, each with
    \code{grow=1}, separated by 1-cell gaps. Resize your terminal
    --- watch the boxes split the new width.
  \item Read \file{maya/src/layout/yoga.cpp}. Identify the three
    phases (measure, distribute, place). Note where the available
    width is subdivided.
\end{enumerate}
